%! Principles of Experimental Design — Unit 1, Lesson 1.6 — Guided Notes & Practice
% THE in-class document, in the MAIN IDEAS / NOTES shape (LESSON_SHAPE.md §1, ported from AP
% Statistics 2026-09-12): the vocabulary box, the hook, then ONE two-column guidednotes table.
% Each row is one idea — a label on the left; on the right one or two COMPLETE printed
% sentences, ONE large pre-drawn display the student reads or names, then a bold prompt and a
% two-across grid of numbered problems with 1.2–2.0cm of writing room. Problems run 1..17.
% DENSITY: a \blank{} only where a word IS the answer (the four cells of the Cedar Ridge table
% in row 2 — the one \wordbank strip sits above it); no mid-sentence blanks; every display is
% pre-drawn; the page belongs to the student's pen.
%   I DO   : the teacher reads the definition, marks up the display, works the first problem
%            of each grid (1, 6, 8, 10).
%   WE DO  : the class works the rest of each grid, students holding the pen; the last row,
%            Guided Practice (14–17), is worked entirely together on Harbor Point instead of
%            Riverbend.
%   YOU DO : nothing on this page — ../homework/ IS the individual practice, started in the
%            last ten minutes of the period. No independent set, no reflection box, no
%            objectives box (the cover carries the targets).
% THE TRAP is problem 4 (a new app against the old packet: an experiment with no "nothing"
% group — the packet group IS the control group). THE CRUX is problem 13 (row 4, "Caused"): a
% bigger self-selected gap is NOT stronger evidence — the design buys the verb, not the gap
% size (PS.DC.3b). The hook vote is re-taken at problem 12. Problem 17 is the crux transferred
% to Harbor Point (25 watched points vs 20 assigned points).
%
% THE DATA (all verified in Python)
%   Riverbend HS experiment: 120 volunteers, generator -> 60 Study Center twice a week / 60 not;
%     45/60 = 75% vs 33/60 = 55%, gap 20 points. Lesson 1.5's watched study on the same Study
%     Center: users 65% vs non-users 80%, gap 15 points the other way.
%   Blocking on last quarter: 80 passed everything -> 40/40; 40 did not -> 20/20; 60 per group.
%   Cedar Ridge Apartments (row 2 fill table): 100 households, 50 reminder / 50 nothing.
%   Guided Practice, Harbor Point pool: 90 volunteers -> 45 at 6 a.m. / 45 at 7 p.m.;
%     27/45 = 60% vs 18/45 = 40%, gap 20; Lesson 1.5 watched 70% vs 45%, gap 25, members chose
%     and 75% of the morning swimmers were retired.
%
% PAGE LOCKSTEP with notes_key/: the key is GENERATED from this file by templates/lesson/mkkey.py
% (spec: templates/lesson/examples/lesson06_notes_key.json) — every \blank{} becomes an \ans{},
% every \vterm{} a \vtermans{}{}, every \par\writelines{n} n \ansline{} calls, and the answer
% argument of every \pcell, \writespace and \labelbox is filled. Answer spaces are fixed
% heights, so the two files cannot drift. KEEP EVERY \pcell ON ONE SOURCE LINE — mkkey matches
% line by line.
% TABLE MECHANICS: inside a cell, `\\` ENDS THE ROW — break lines with \par. A row cannot break
% across pages: the figure gets its own sub-row (`\\` then `& ...`) and EACH GRID ROW is its own
% sub-row — close the probgrid, `\\ &`, reopen as probgrid* (no top rule).
\documentclass[12pt]{article}
\usepackage{saar-article}
\usepackage{saar-boxes}
\newcommand{\TallMath}[1]{$\displaystyle #1\rule[-1.4em]{0pt}{3.2em}$}
% Word bank strip — ONLY above a table the student fills (row 2). Defined per document;
% the key inherits it and prints the same strip, so heights match.
\newcommand{\wordbank}[1]{\par\vspace{2pt}\noindent\fcolorbox{goldacc}{goldbg}{\parbox{\dimexpr\linewidth-2\fboxsep-2\fboxrule\relax}{\small\textbf{Word bank:} #1}}\par\vspace{4pt}}
% Vocabulary rows: a FIXED-HEIGHT row carrying the term label and OPEN writing space —
% no underline beside the term and no rule line (user correction 2026-09-06: the black
% answer line crowded the row; students write beside and below the term). The key fills
% exactly the same height, so the box cannot drift blank vs. keyed. saar-article's
% \termblank adds an inline blank and a rule line, so the pair is defined here (the key is
% generated from this file and inherits both). 2.15cm per row gives students three
% handwritten lines of room (trimmed from 1.2's 2.3cm so the table's first row still
% starts on page 1); keep key definitions to two lines.
\newcommand{\termrowheight}{2.15cm}
\newlength{\termrowinset}\setlength{\termrowinset}{7pt}
\newcommand{\vterm}[1]{%
  \par\noindent\begin{minipage}[t][\termrowheight][t]{\linewidth}%
    \vspace*{\termrowinset}%
    \textbf{\textcolor{cerulean}{#1:}}%
  \end{minipage}\par}
\newcommand{\vtermans}[2]{%
  \par\noindent\begin{minipage}[t][\termrowheight][t]{\linewidth}%
    \vspace*{\termrowinset}%
    \textbf{\textcolor{cerulean}{#1:}}\quad\ans{#2}%
  \end{minipage}\par}

\begin{document}

\pageheader{Unit 1, Lesson 1.6}{Guided Notes \& Practice}
% NAMESTRIP: no \namedateperiod here. The student writes their name once, on the
% cover this component is stapled behind.

\vspace{0.08in}

\begin{vocabbox}
\small
Fill in each term \textbf{as we name it} in the notes below.
\vterm{Experiment}
\vterm{Control group}
\vterm{Random assignment}
\vterm{Confounding variable}
\vterm{Blinding}
\end{vocabbox}

\boxguard[12]
\begin{hookbox}
\small
Last quarter the Riverbend counselor \emph{watched}: $65\%$ of Study Center users passed
every class, against $80\%$ of those who never went --- it looked \emph{harmful}. This
quarter she \emph{assigned}: $120$ volunteers split by a generator gave $75\%$ against
$55\%$ --- now it looks \emph{helpful}. \textbf{Same school, same Study Center, opposite
conclusions --- which study should she believe?} Circle one: \; last quarter's \quad this
quarter's \quad neither \; --- and say why. We vote again at problem~12.
\par\writelines{2}
\end{hookbox}

\small
\begin{guidednotes}
% ── ROW 1: assigning, not watching ───────────────────────────────────────────
\mainidea[Assign, don't watch]{Experiment} &
In an \textbf{experiment} the researcher \textbf{assigns} the individuals to groups and
\textbf{imposes a treatment} on at least one. The individuals are the \textbf{experimental
units} (\textbf{subjects}, when they are people); what is applied is the \textbf{treatment};
the group that does \emph{not} get it is the \textbf{control group} --- the baseline, not a
wasted group. \textbf{The test: who built the groups?}
\\
& {\centering
\begin{tikzpicture}[scale=0.66]
  % the 120 volunteers
  \fill[frost, rounded corners=5pt] (0,0.9) rectangle (2.6,2.3);
  \draw[cerulean, thick, rounded corners=5pt] (0,0.9) rectangle (2.6,2.3);
  \node[font=\scriptsize\bfseries\color{cerulean}, align=center] at (1.3,1.6) {$120$\\volunteers};
  % the generator
  \draw[->, very thick, charcoal] (2.7,1.6) -- (3.7,1.6);
  \fill[goldbg, rounded corners=5pt] (3.8,0.9) rectangle (6.6,2.3);
  \draw[goldacc, very thick, rounded corners=5pt] (3.8,0.9) rectangle (6.6,2.3);
  \node[font=\scriptsize\bfseries\color{goldacc}, align=center] at (5.2,1.6) {random number\\generator};
  % the two groups
  \draw[->, very thick, charcoal] (6.7,1.9) -- (7.7,2.7);
  \draw[->, very thick, charcoal] (6.7,1.3) -- (7.7,0.5);
  \fill[goldbg, rounded corners=5pt] (7.8,2.15) rectangle (11.4,3.25);
  \draw[goldacc, thick, rounded corners=5pt] (7.8,2.15) rectangle (11.4,3.25);
  \node[font=\scriptsize\bfseries, align=center] at (9.6,2.7) {$60$ --- Study Center\\twice a week};
  \fill[frost, rounded corners=5pt] (7.8,-0.05) rectangle (11.4,1.05);
  \draw[cerulean, thick, rounded corners=5pt] (7.8,-0.05) rectangle (11.4,1.05);
  \node[font=\scriptsize\bfseries, align=center] at (9.6,0.5) {$60$ --- no\\Study Center};
  % the results
  \draw[->, thick, charcoal] (11.5,2.7) -- (12.3,2.7);
  \draw[->, thick, charcoal] (11.5,0.5) -- (12.3,0.5);
  \node[font=\scriptsize, align=left, anchor=west] at (12.35,2.7) {$45$ passed\\every class: $75\%$};
  \node[font=\scriptsize, align=left, anchor=west] at (12.35,0.5) {$33$ passed\\every class: $55\%$};
  \node[font=\tiny\itshape\color{slate}] at (5.2,0.45) {nobody chose --- chance did};
\end{tikzpicture}\par\vspace{5pt}
The $120$ are the \labelbox{2cm}{}\quad attending is the \labelbox{2cm}{}\quad the other $60$ are the \labelbox{2.4cm}{}\par}
\\
& \notesprompt{Experiment or not? Say who decided which group each person was in.}
\begin{probgrid}{|Y|Y|}
\pcell{1}{Pull $150$ student records and compare the pass rate of students who used the Study Center to those who did not.}{1.2cm}{} & \pcell{2}{Ask for volunteers, then let a random number generator send $60$ to the Study Center and $60$ nowhere.}{1.2cm}{} \\ \hline
\end{probgrid}
\\
& \begin{probgrid}*{|Y|Y|}
\pcell{3}{Survey a random sample of pool members about what time they swim and how well they sleep.}{1.2cm}{} & \pcell{4}{Give $60$ randomly chosen volunteers a new tutoring app and the other $60$ the old worksheet packet. \textbf{Vote first:} is there a control group?}{1.5cm}{} \\ \hline
\end{probgrid}
\\ \hline
% ── ROW 2: the four principles ───────────────────────────────────────────────
\mainidea[Four]{Principles} &
Almost every experiment worth trusting does all four of these. A variable that differs
between the groups \emph{and} also affects the response is a \textbf{confounding variable}
--- Lesson 1.5's lurking variable, now \emph{inside} the experiment. \textbf{Randomization
makes the groups alike before you start; control keeps them alike while you run.}
\\
& {\centering
\begin{tikzpicture}[scale=0.85]
  \foreach \i/\cname/\cdesc in {
      0/Comparison/{at least two groups,\\one of them a\\control group},
      1/Randomization/{chance decides the\\groups --- not the\\researcher, not the\\subject, not the\\sign-up order},
      2/Replication/{many units in\\each group, so one\\person's luck\\cannot decide it},
      3/Control/{every other\\condition the same\\for both groups}} {
    \fill[frost, rounded corners=4pt] ({\i*2.9},0) rectangle ({\i*2.9+2.7},2.2);
    \draw[cerulean, thick, rounded corners=4pt] ({\i*2.9},0) rectangle ({\i*2.9+2.7},2.2);
    \node[font=\scriptsize\bfseries\color{cerulean}] at ({\i*2.9+1.35},1.9) {\cname};
    \node[font=\tiny, align=center] at ({\i*2.9+1.35},0.85) {\cdesc};
  }
\end{tikzpicture}\par}
\\
& \textbf{5.} \emph{Cedar Ridge Apartments:} $100$ households volunteer; the manager randomly
assigns $50$ a weekly recycling reminder and $50$ nothing, then weighs their recycling for a
month. Each row changes the plan --- name the principle it breaks.
\wordbank{comparison \;\textbullet\; randomization \;\textbullet\; replication \;\textbullet\; control}
\footnotesize\renewcommand{\arraystretch}{1.0}
\begin{tabularx}{\linewidth}{@{} X p{2.6cm} @{}}
\toprule
\textbf{Instead, the manager\ldots} & \textbf{Principle broken} \\
\midrule
sends the reminder to all $100$ households.
  & \blank{2.2cm} \\
sends it to the first $50$ households that signed up.
  & \blank{2.2cm} \\
sends it to $1$ household and nothing to $1$ other household.
  & \blank{2.2cm} \\
sends it to the $50$ in the renovated building, which just got new bins.
  & \blank{2.2cm} \\
\bottomrule
\end{tabularx}\small
\\
& \notesprompt{Now check Riverbend's design.}
\begin{probgrid}{|Y|Y|}
\pcell{6}{One short phrase each: how does Riverbend's experiment meet comparison, randomization, replication, and control?}{1.5cm}{} & \pcell{7}{In Cedar Ridge's last plan the reminder group also got new bins. Name the confounding variable, and say why the manager cannot tell what caused the extra recycling.}{1.5cm}{} \\ \hline
\end{probgrid}
\\ \hline
% ── ROW 3: blinding and blocking ─────────────────────────────────────────────
\mainidea[Two refinements]{Blinding \& Blocking} &
A subject who responds to the \emph{idea} of being treated shows the \textbf{placebo
effect}; a \textbf{placebo} is a fake treatment built to look real. \textbf{Blinding} hides
who is in which group from the subjects (\textbf{single-blind}), or from them \emph{and} the
people scoring the response (\textbf{double-blind}). A \textbf{randomized block design}
sorts the units into \textbf{blocks} of similar units first, then randomizes \emph{inside}
each block (\textbf{matched pairs} is the block of two); Riverbend's plain $60/60$ split is a
\textbf{completely randomized design}.
\\
& {\centering
\begin{tikzpicture}[scale=0.74]
  % left panel: blinding
  \fill[goldbg, rounded corners=5pt] (0,0) rectangle (5.8,3.2);
  \draw[goldacc, thick, rounded corners=5pt] (0,0) rectangle (5.8,3.2);
  \node[font=\footnotesize\bfseries\color{goldacc}] at (2.9,2.85) {BLINDING};
  \draw[charcoal, thick] (0.4,1.45) rectangle (2.8,2.45);
  \node[font=\scriptsize\bfseries, align=center] at (1.6,1.95) {Room A\\tutoring};
  \draw[charcoal, thick] (3.0,1.45) rectangle (5.4,2.45);
  \node[font=\scriptsize\bfseries, align=center] at (4.2,1.95) {Room B\\quiet study};
  \node[font=\tiny\itshape, align=center] at (2.9,0.95) {students: not told which room is the program\\graders: not told who sat where};
  \node[font=\tiny\bfseries\color{goldacc}] at (2.9,0.3) {nobody knows who was picked};
  % right panel: blocking
  \fill[frost, rounded corners=5pt] (6.6,0) rectangle (13.0,3.2);
  \draw[cerulean, thick, rounded corners=5pt] (6.6,0) rectangle (13.0,3.2);
  \node[font=\footnotesize\bfseries\color{cerulean}] at (9.8,2.85) {BLOCKING ON GRADES};
  \node[font=\scriptsize\bfseries] at (7.1,1.55) {$120$};
  \draw[->, thick, charcoal] (7.4,1.75) -- (7.9,2.15);
  \draw[->, thick, charcoal] (7.4,1.35) -- (7.9,0.95);
  \draw[charcoal, thick] (8.0,1.85) rectangle (10.3,2.55);
  \node[font=\scriptsize, align=center] at (9.15,2.2) {passed all:\\$80$};
  \draw[charcoal, thick] (8.0,0.6) rectangle (10.3,1.3);
  \node[font=\scriptsize, align=center] at (9.15,0.95) {did not:\\$40$};
  \draw[->, thick, charcoal] (10.4,2.2) -- (10.9,2.2);
  \draw[->, thick, charcoal] (10.4,0.95) -- (10.9,0.95);
  \node[font=\tiny, align=left, anchor=west] at (10.9,2.2) {$?$ Center\\$?$ control};
  \node[font=\tiny, align=left, anchor=west] at (10.9,0.95) {$?$ Center\\$?$ control};
  \node[font=\tiny\itshape\color{cerulean}] at (9.8,0.25) {chance splits \emph{inside} each block};
\end{tikzpicture}\par\vspace{4pt}
Room B is a \labelbox{2cm}{}\quad neither side knows: \labelbox{2.6cm}{}\par}
\\
& \notesprompt{Who knows, and who is in which block?}
\begin{probgrid}{|Y|Y|}
\pcell{8}{Students picked for the Study Center try harder simply \emph{because} they were picked. What is that called, and how does Room~B remove it?}{1.5cm}{} & \pcell{9}{Fill the blocking diagram: of the $80$ who passed everything, how many go to each group? Of the $40$ who did not? Why block instead of trusting chance?}{1.5cm}{} \\ \hline
\end{probgrid}
\\ \hline
% ── ROW 4: THE CRUX — what buys the verb ─────────────────────────────────────
\mainidea[What buys the verb]{Caused} &
Watching earns an \textbf{association}. Assigning earns more: because chance built the two
groups, they were alike \emph{before} the treatment, so the treatment is the only thing left
that can explain the gap. \textbf{That is what buys the word \emph{caused} --- the design,
not the size of the gap.}
\\
& {\centering
\begin{tikzpicture}[scale=0.78]
  % left panel: the watched study
  \fill[frost, rounded corners=5pt] (0,0) rectangle (5.7,3.0);
  \draw[cerulean, thick, rounded corners=5pt] (0,0) rectangle (5.7,3.0);
  \node[font=\scriptsize\bfseries\color{cerulean}] at (2.85,2.7) {LAST QUARTER: WATCHED};
  \node[font=\tiny, align=center] at (2.85,2.15) {$150$ records; the students \textbf{chose}\\whether to use the Study Center};
  \node[font=\small\bfseries, align=center] at (2.85,1.35) {users $65\%$\\vs.\ non-users $80\%$};
  \node[font=\tiny\itshape] at (2.85,0.62) {gap: $15$ points --- the wrong way};
  \node[font=\tiny\itshape\color{slate}] at (2.85,0.27) {who chose? the already struggling};
  % right panel: the experiment
  \fill[goldbg, rounded corners=5pt] (6.5,0) rectangle (12.2,3.0);
  \draw[goldacc, thick, rounded corners=5pt] (6.5,0) rectangle (12.2,3.0);
  \node[font=\scriptsize\bfseries\color{goldacc}] at (9.35,2.7) {THIS QUARTER: ASSIGNED};
  \node[font=\tiny, align=center] at (9.35,2.15) {$120$ volunteers; a \textbf{generator} chose\\who used the Study Center};
  \node[font=\small\bfseries, align=center] at (9.35,1.35) {treated $75\%$\\vs.\ control $55\%$};
  \node[font=\tiny\itshape] at (9.35,0.62) {gap: $20$ points};
  \node[font=\tiny\itshape\color{slate}] at (9.35,0.27) {who chose? nobody --- groups started alike};
\end{tikzpicture}\par\vspace{4pt}
Last quarter may claim: \labelbox{2.6cm}{}\qquad This quarter may claim: \labelbox{2.2cm}{}\par}
\\
& \textbf{In your own words:} why may this quarter say \emph{caused}?
\writespace{1.0cm}{}
\\
& \notesprompt{Now settle the hook vote.}
\begin{probgrid}{|Y|Y|}
\pcell{10}{Last quarter's Study Center users were mostly students who were \emph{already struggling}. Where did the struggling volunteers end up this quarter --- and what does that do to the gap?}{1.3cm}{} & \pcell{11}{Write this quarter's conclusion in one sentence for the counselor. Use the verb the design earns.}{1.3cm}{} \\ \hline
\end{probgrid}
\\
& \begin{probgrid}*{|Y|Y|}
\pcell{12}{\textbf{Vote again:} last quarter's \; this quarter's \; neither. Give the reason. A classmate says ``this quarter's, because $20$ points beats $15$.'' Is that the reason?}{1.4cm}{} & \pcell{13}{Suppose last quarter's \emph{watched} study had found $90\%$ vs.\ $55\%$ --- a $35$-point gap, bigger than the experiment's $20$. Would that be the study to believe? Explain.}{1.4cm}{} \\ \hline
\end{probgrid}
\\ \hline
% ── ROW 5: GUIDED PRACTICE — we do, together, a fresh context ────────────────
\mainidea[Guided practice]{Harbor Point Swim Times} &
Last lesson the pool \emph{watched}: $70\%$ of morning swimmers slept $7+$ hours against
$45\%$ of evening swimmers --- but members chose their times, and $75\%$ of the morning
swimmers were retired. Now the director \emph{assigns}: $90$ volunteers, a generator, $45$ to
a reserved $6$~a.m.\ slot and $45$ to $7$~p.m.\ for four weeks. Each keeps a sleep log.
\\
& {\centering
\begin{tikzpicture}[scale=0.66, every node/.append style={font=\tiny}]
  \fill[frost, rounded corners=5pt] (0,0.9) rectangle (2.6,2.3);
  \draw[cerulean, thick, rounded corners=5pt] (0,0.9) rectangle (2.6,2.3);
  \node[font=\scriptsize\bfseries\color{cerulean}, align=center] at (1.3,1.6) {$90$\\volunteers};
  \draw[->, very thick, charcoal] (2.7,1.6) -- (3.7,1.6);
  \fill[goldbg, rounded corners=5pt] (3.8,0.9) rectangle (6.6,2.3);
  \draw[goldacc, very thick, rounded corners=5pt] (3.8,0.9) rectangle (6.6,2.3);
  \node[font=\scriptsize\bfseries\color{goldacc}, align=center] at (5.2,1.6) {random number\\generator};
  \draw[->, very thick, charcoal] (6.7,1.9) -- (7.7,2.7);
  \draw[->, very thick, charcoal] (6.7,1.3) -- (7.7,0.5);
  \fill[goldbg, rounded corners=5pt] (7.8,2.15) rectangle (11.4,3.25);
  \draw[goldacc, thick, rounded corners=5pt] (7.8,2.15) rectangle (11.4,3.25);
  \node[font=\scriptsize\bfseries, align=center] at (9.6,2.7) {$45$ --- swim at\\$6$~a.m.};
  \fill[frost, rounded corners=5pt] (7.8,-0.05) rectangle (11.4,1.05);
  \draw[cerulean, thick, rounded corners=5pt] (7.8,-0.05) rectangle (11.4,1.05);
  \node[font=\scriptsize\bfseries, align=center] at (9.6,0.5) {$45$ --- swim at\\$7$~p.m.};
  \draw[->, thick, charcoal] (11.5,2.7) -- (12.3,2.7);
  \draw[->, thick, charcoal] (11.5,0.5) -- (12.3,0.5);
  \node[font=\scriptsize, align=left, anchor=west] at (12.35,2.7) {$27$ averaged\\$7+$ hours};
  \node[font=\scriptsize, align=left, anchor=west] at (12.35,0.5) {$18$ averaged\\$7+$ hours};
  \node[font=\tiny\itshape\color{slate}] at (5.2,0.45) {last lesson, watching: $70\%$ vs.\ $45\%$};
\end{tikzpicture}\par}
\\
& \notesprompt{We work these together.}
\begin{probgrid}{|Y|Y|}
\pcell{14}{Name the experimental units, the treatment, and the control group. Who decided the slots?}{1.4cm}{} & \pcell{15}{What percent of each group averaged $7+$ hours? How many points apart?}{1.4cm}{} \\ \hline
\end{probgrid}
\\
& \begin{probgrid}*{|Y|Y|}
\pcell{16}{One phrase each: how does this design meet the four principles? Who can be blinded here?}{1.3cm}{} & \pcell{17}{Watching gave $25$ points; assigning gives $20$. Which is the stronger evidence, and why?}{1.3cm}{} \\ \hline
\end{probgrid}
\\ \hline
\end{guidednotes}

% ── THE NOTES END AT GUIDED PRACTICE ─────────────────────────────────────────
% There is deliberately no "you do" here: ../homework/ is the individual practice,
% started in class in the last ten minutes while the teacher circulates.

\end{document}
